\pdfoutput=1
\documentclass[a4paper]{base}
\usepackage{chem}

\newcommand{\pg}[1][1]{\M{\ifthenelse{\equal{#1}{1}}{\psi^{\prime}}{\psi^{(#1)}}}}
\newcommand{\stoch}[1]{\M{\langle #1 \rangle}}
\newcommand{\crit}[1]{\M{#1^{*}}}
\newcommand{\cib}[0]{\M{\beta}}
\newcommand{\pdn}[1]{\partial_{\,\dn{#1}}}


\begin{document}
    \title{Parametrizing A Population Balance Equation}
    \author[R.A. Milton]{Robert A. Milton}
    \email{r.a.milton@sheffield.ac.uk}
    \address{Department of Chemical, Materials and Biological Engineering, The University of Sheffield, Sheffield~S1~3JD, United Kingdom}
    \thanks{This work was funded by Solomon f. Brown.}
    \subjclass[2010]{}
    \keywords{Asymptotics}
    \begin{abstract}
        Surely this cannot be that difficult.
    \end{abstract}
    \maketitle


\section{Introduction}\label{sec:intro}
    A plethora of quantities arise in this work, so we endeavour to consistently indicate their nature at a glance.
    All constants $\norm{\cdot},\stoch{\cdot},\mds{\cdot}\in\overline{\mathbb{R}^{+}}$ are non-negative, and enclosed in \norm{\cdot} or \stoch{\cdot} without meaning (there are no norms or inner products here), or \mds{\cdot} only when taking the absolute value of $\cdot$.
    All \M{italic} quantities are dimensionless, to ease the unavoidable chore of tracking units.
    Undefined notation adheres to the Digital Library of Mathematical Functions \cite{NIST:DLMF}.
    
    
    A number $\dn{n(t,L)}~\un{m^{-4}}$ of particles of length L~\un{m} grow within each unit of reactor volume V~\un{m^{3}} as seconds t~\un{s} pass. 
    Conserving mass, the Population Balance Equation is
    \begin{equation}\label{eq:Intro:PBE}
        \begin{aligned}
            \dn{\partial_{t} (V n(t,L))} &= -\dn{G^{+} \pdn{L} (V \dn{n}(t,L))}\\
            \dn{n(0, L)} &= 0 \\
            \dn{n(t,0)} &= \dn{B^{+}/\; G^{+}}
        \end{aligned}
    \end{equation}
    where $\dn{B^{+}}$ \un{m^{-3}s^{-1}} and $\dn{G^{+}}$ \un{m\,s^{-1}} are particle birth (nucleation) and growth rates. These relate empirically to supersaturation $S>1$:
    \begin{equation}\label{def:Intro:BG+}
        \begin{aligned}
            \QT{birth} \dn{B^{+}} &\deq \norm{\dn{B}} \left(S-1\right)^{+\mds{b}} \\
            \QT{growth} \dn{G^{+}} &\deq \norm{\dn{G}} \left(S-1\right)^{+\mds{g}}
        \end{aligned}
    \end{equation}
    where $(\#)^{+} \deq \max(0, \#)$ denotes the positive part of $(\#)$. In experiments, critical supersaturation $S=1$ might occur many times, but we cannot allow this. We must insist that
    \begin{equation}\label{def:Intro:crit}
        \begin{aligned}
            \QT{supersaturation} S \in \opclo{1}{\sup S} \QT{at all times} \dn{t}>\crit{\dn{t}}
        \end{aligned}
    \end{equation}
    for finite supersaturation $\sup{S}\in\bb{R}$ and critical time $\crit{\dn{t}}\in\bb{R}$ at which $S$ first exceeds 1. This is not limiting per se, but it does mean that each critical instant in an experiment theoretically begins a new experiment. 
    Throughout this work, an asterisk always denotes the critical instant $t=S-1=0$

    Given lengthscale $\norm{\dn{L}}\un{m}$ and timescale $\norm{\dn{t}}\un{s}$, dimensionless prefactors are defined as
    \begin{equation}\label{def:Intro:BG}
        \begin{aligned}
            \QT{birth} \norm{B} &\deq \norm{\dn{B}} \norm{\dn{L}}^{3}\norm{\dn{t}} \\
            \QT{growth} \norm{G} &\deq \norm{\dn{G}} \norm{\dn{L}}^{-1}\norm{\dn{t}}
        \end{aligned}
    \end{equation}
    together with dimensionless
    \begin{equation}\label{def:Intro:tV}
        \begin{aligned}
            \QT{time}\qquad t &\deq \norm{\dn{t}}^{-1} (\dn{t}-\crit{\dn{t}}) \\
            \QT{volume} V(t) &\deq \norm{\dn{L}}^{-3} \dn{V}
        \end{aligned}
    \end{equation}
    This work aims to elucidate the dimensionless empirical birth and growth parameters $\norm{B},\mds{b},\norm{G},\mds{g}$.
    

\section{The Crystal Size Distribution (CSD)}\label{sec:CSD}
    The Population Balance Equation conveys a symmetry between synthetic `time' $T$ and `length' $L$ defined as
    \begin{subequations} \label{def:CSD:TL}
        \begin{align}
            S(t) &\deq \max (S,1) \\
            T &\deq \norm{\dn{L}}^{-1} \dn{\int_{0}^{t} G^{+} \, d{t}} 
                = \int_{0}^{t} \norm{G} \left(S(\dot{t})-1\right)^{\mds{g}} \!\dd{\dot{t}} \label{def:CSD:T} \\[3pt]
            L &\deq T -  \norm{\dn{L}}^{-1} \dn{L}
        \end{align}
    \end{subequations}
    These forms are chosen to render the PBE linear and invariant to exchanging $T$ and $L$.
    This symmetry manifests in the particle count $N(T)f(T,L)$ defined as
    \begin{subequations} \label{def:CSD:Nf}
        \begin{align}
            N(\bb{R^{-}}) &\deq \set{0} \\
            N(T) &\deq \dn{\int_{0}^{t} B^{+} V \, d\dot{t}} 
                = \int_{0}^{t} \norm{B} \left(S(\dot{t})-1\right)^{\mds{b}} V(\dot{t}) \dd{\dot{t}} \label{def:CSD:N} \\[2pt]
            f(T,L) &\deq N(T)^{-1} \norm{\dn{L}}^{-1} \dn{V n(t,L)}
        \end{align}
    \end{subequations}

    Substituting these variables in Population Balance \Cref{eq:Intro:PBE} immediately yields
    \begin{equation} \label{sol:CSD:PBE}
        \begin{aligned}
            \partial_T N(T) f(T,L) &= 0 \\
            N(T) f(T, L) &= \partial_{t} N(T)/ \partial_{t} T = \partial_T N(T)
        \end{aligned}
    \end{equation}
    These two lines combine to establish
    \begin{equation} \label{sol:CSD:N}
        \begin{aligned}
            f(T, L) &= N(T)^{-1} \partial_L N(L) \\
            N(T)^{-1} N(L)  &= \int_{0}^{L} f(T, \dot{L}) \dd{\dot{L}}
        \end{aligned}
    \end{equation}

    The Crystal Size Distribution (CSD) at `time' $T$ is $N(T)^{-1} N(L)$, its probability density $f(T, L)$, and its moments
    \begin{equation} \label{def:CSD:M}
        \begin{aligned}
             M_0(T) &\deq \int_{0}^{T} f(T, L)(T-L)^{0} \dd{L} \;= N(T)^{-1}N(T) = 1 \\
             M_m(T) &\deq \int_{0}^{T} f(T, L)(T-L)^{m} \dd{L} = N(T)^{-1} m \int_{0}^{T} N(L)M_{m-1}(L) \dd{L}
        \end{aligned}
    \end{equation}
    The number of particles born by `time' $T$ is $N(T)$, and the number of particles whose `length' does not exceed $L$ is $N(L)$.
    This is the symmetry conveyed by the PBE: `time' $T$ and `length' $L$ are effectively the same thing.
    Like a growing child, a crystal equates how old it is with how big it is.

    The CSD $N(T)^{-1} N(L)$ is the fraction of particles at `time' $T$ whose `length' does not exceed $L$, which is the same as the fraction of particles at `time' $T$ born by `time' $L$.

\section{Taylor Approximation}\label{sec:Taylor}
    This Section determines the CSD from supersaturation $S(t)$ and volume $V(t)$, expanded about $t=\bar{t}$ as
    \begin{equation} \label{def:Taylor:s}
        \begin{aligned}
            \ln (S(t)-1) &\phantom{:}= \sum_{k=0}^{K} \frac{\bar{\sigma}^{(k)}}{\Gamma(k+1)} \left(t-\bar{t}\right)^{k}
                +  O \left(\frac{\bar{\sigma}^{(K+1)}}{\Gamma(K+2)}\left(t-\bar{t}\right)^{K+1}\right)\\
            \QT{where}\bar{\sigma}^{(k)} &\deq \partial_{\bar{t}}^{\,k} \ln (S(\bar{t})-1) \\
            V(t) &\phantom{:}= \sum_{k=0}^{K} \frac{\bar{V}^{(k)}}{\Gamma(k+1)} \left(t-\bar{t}\right)^{k}
                +  O \left(\bar{V}^{(K+1)} \frac{\left(t-\bar{t}\right)^{K+1}}{\Gamma(K+2)}\right)\\
            \QT{where}\bar{V}^{(k)} &\deq \partial_{\bar{t}}^{\,k} V(\bar{t})
        \end{aligned}
    \end{equation}
    The partial derivatives here are normally provided by an exogenous model of supersaturation and mixing.


\subsection{Linear Approximation}\label{sec:Lin}
    Taking $K=1$ requires
    \begin{equation} \label{sol:Quad:N}
        \begin{aligned}
           \bar{\sigma}^{\prime} > 0
        \end{aligned}
    \end{equation}
    Integrating \Cref{def:CSD:N} using \Cref{def:Taylor:s} with $K=1$ yields
    \begin{equation} \label{sol:Lin:N}
        \begin{aligned}
            \norm{N}^{-1} N(T) &\phantom{:}= \left(\bar{v} + \bar{V}^{\prime}t\right) \exp \left(\norm{b} t\right) - \bar{v} \\
            \QT{where}\norm{b} &\deq \mds{b}\bar{\sigma}^{\prime} \\
            \norm{\bar{\sigma}} &\deq \mds{b}\bar{\sigma}-\norm{b} \bar{t} \\
            \norm{N} &\deq \norm{B} \norm{b}^{-1} \exp \norm{\bar{\sigma}} \\
            \bar{v} &\deq \bar{V} - \bar{V}^{\prime}\bar{t} - \bar{V}^{\prime} \norm{b}^{-1}
        \end{aligned}
    \end{equation}
    Integrating \Cref{def:CSD:T} for small $\mds{g}$ using \Cref{def:Taylor:s} with $K=1$ yields
    \begin{equation} \label{sol:Lin:t}
        \begin{aligned}
            t &= \frac{T}{\norm{G}} - \frac{\mds{g}}{\mds{b}}\frac{T}{\norm{G}} \left(\frac{\norm{b}}{2}\frac{T}{\norm{G}} + \norm{\bar{\sigma}}\right) 
                + O(\mds{g}^{2})
        \end{aligned}
    \end{equation}
    so
    \begin{equation} \label{sol:Lin:CSD}
        \begin{aligned}
            \norm{N}^{-1} N(T) &= \left(\bar{v} + \bar{V}^{\prime}\frac{T}{\norm{G}}\right) \exp \left(\norm{b} \frac{T}{\norm{G}}\right) - \bar{v} \\
            &\phantom{=\ }-\frac{\mds{g}}{\mds{b}}\frac{T}{\norm{G}} \left(\frac{\norm{b}}{2}\frac{T}{\norm{G}} + \norm{\bar{\sigma}}\right)
                 \exp \left(\norm{b} \frac{T}{\norm{G}}\right)
                 \left(\norm{b}\left(\bar{v} + \bar{V}^{\prime}\frac{T}{\norm{G}}\right) + \bar{V}^{\prime}\right) \\
            &\phantom{=\ }+ O(\mds{g}^{2})
        \end{aligned}
    \end{equation}

\subsection{Quadratic Approximation}\label{sec:Quad}
    Taking $K=2$, integration ``completes the square'', effectively shifting $\bar{t}$ to the maximum of $\ln (S{t}-1)$ so that
    \begin{equation} \label{sol:Quad:N}
        \begin{aligned}
           \bar{\sigma}^{\prime} \deq 0 > \bar{\sigma}^{\prime\prime}
        \end{aligned}
    \end{equation}
    Integrating \Cref{def:CSD:N} using \Cref{def:Taylor:s} with $K=2$ yields
    \begin{equation} \label{sol:Quad:N}
        \begin{aligned}
            \norm{N}^{-1} N(T) &\phantom{:}= \bar{v} \left(\erf(\norm{b}(t-\bar{t}))+\erf(\norm{b}\bar{t})\right)  
             - \left(\frac{\bar{V}^{\prime\prime}}{2}t + \bar{v}^{\prime}\right) \exp\left(\norm{b}^{2}t(2\bar{t}-t)\right)
            + \bar{v}^{\prime} \\
            \QT{where}\norm{b} &\deq \mds{\frac{\mds{b}\bar{\sigma}^{\prime\prime}}{2}}^{1/2} \\
            \norm{\bar{\sigma}} &\deq \mds{b}\bar{\sigma}-\norm{b}^2 \bar{t}^2 \\
            \norm{N} &\deq \frac{\norm{B}\exp \norm{\bar{\sigma}}} {2\norm{b}^{2}}  \\
            \bar{v} &\deq \norm{b} \pi^{1/2} \exp \left(\norm{b}^{2}\bar{t}^{2}\right) \left(\bar{V} + \frac{\bar{V}^{\prime\prime}}{4\norm{b}^{2}}\right) \\
            \bar{v}^{\prime} &\deq \bar{V}^{\prime} - \frac{\bar{V}^{\prime\prime}\bar{t}}{2}\\
        \end{aligned}
    \end{equation}
    Integrating \Cref{def:CSD:T} for small $\mds{g}$ using \Cref{def:Taylor:s} with $K=1$ yields
    \begin{equation} \label{sol:Quad:t}
        \begin{aligned}
            t &= \frac{T}{\norm{G}} - \frac{\mds{g}}{\mds{b}}\frac{T}{\norm{G}} 
            \left(\norm{b}^{2} \frac{T}{\norm{G}} \left(\bar{t}-\frac{1}{3}\frac{T}{\norm{G}}\right) + \norm{\bar{\sigma}}\right) + O(\mds{g}^{2})
        \end{aligned}
    \end{equation}


%%%%%%%%%%%%%%%%%%%%%%%%%%%%%%%%%%%%%%%%%%%%%%%%%%%%%%%%%%%%%%%%%%%%%%%%%%%%%%%%%%%%%%%%%%%%%%%%%%%%%%%%%%%%%%%%%%%

\bibliographystyle{amsplain} 
\bibliography{../References}
\end{document}
\endinput
